\chapter{Preguntas de portafolio: pruebas preregistradas}
\label{chap:anexo-portafolio}

Los dos preregistros de este anexo no extienden la pregunta de la tesis: la aplican. Llevan sus resultados a lo que preguntaría quien gestiona un portafolio de deuda soberana emergente: si la fragilidad o los fundamentos pronostican el spread, si ayudan a cubrirlo y si la simultaneidad ensancha su riesgo. Ninguna de sus seis hipótesis primarias pasa la corrección de Holm. Van aquí, y no en el Anexo~\ref{chap:anexo-extensiones}, para separar lo que amplía la evidencia de la tesis de lo que la pone a prueba como herramienta.

\section{Una señal de valor relativo y la cola de los cambios del spread: pruebas preregistradas}
\label{sec:quant-prereg}

Un séptimo preregistro (\texttt{4\_Redaccion/preregistro\_quant.md}), escrito antes de construir cualquier retorno, lleva dos resultados de la tesis a las preguntas que haría quien gestiona un portafolio de deuda soberana emergente. La primera es si la precedencia mensual de la fragilidad (Sección~\ref{sec:mensual}) sirve para ordenar a los países según su retorno del mes siguiente (Q1). La segunda es si la simultaneidad de los dos riesgos, que no eleva el nivel del spread más que la suma, ensancha la distribución de sus cambios (Q3). Todo es mensual, de enero de 2015 en adelante (139 meses), y sin anticipación: el balance de los bancos se usa tres meses después de su fecha de cierre, el riesgo de cola es el del trimestre $q-2$ y el tercil alto de cada riesgo se define con la historia de cada país disponible hasta el mes anterior. El retorno de exceso se aproxima con una duración de siete años, $r\approx-7\,\Delta s+s/12$, porque no hay índices de retorno total por país en el repositorio, y cada operación paga la mitad de un \textit{bid-ask} de 15 puntos básicos de spread.

\begin{table}[htbp]
\centering
\small
\caption[Señal de valor relativo y cola de los cambios del spread]{Señal de valor relativo con la fragilidad (Q1) y cola de los cambios del spread bajo simultaneidad (Q3): pruebas preregistradas}
\label{tab:quant-prereg}
\begin{adjustbox}{max width=\linewidth}
\begin{tabular}{>{\raggedright\arraybackslash}p{4.6cm}>{\raggedright\arraybackslash}p{7.6cm}>{\raggedright\arraybackslash}p{2.6cm}}
\toprule
Hipótesis & Resultado & Holm \\
\midrule
Q1: cartera larga/corta por terciles del cambio mensual de la distancia al incumplimiento & Sharpe anual neto $-4{,}07$ (IC 90\,\%: $-4{,}93$ a $-3{,}17$); bruto $+0{,}01$ ($-0{,}44$ a $+0{,}43$); Fama--MacBeth, $t=-0{,}47$; rotación de 2,5 veces por mes & $1{,}00$: no pasa \\
\addlinespace
Q3-H1: escala de los cambios mensuales del spread (MMQR) & Fragilidad alta $+0{,}062$ ($p=0{,}005$); riesgo de cola alto $+0{,}019$ ($0{,}085$); su coincidencia $-0{,}046$ ($0{,}97$) & $1{,}00$: no pasa \\
\addlinespace
Q3-H2: VaR al 95\,\% fuera de la muestra & Condicional frente a EWMA: Diebold--Mariano $p=0{,}47$; cobertura condicional rechazada ($p=0{,}002$); el VaR histórico de 60 meses es el mejor de los tres & $1{,}00$: no pasa \\
\bottomrule
\end{tabular}
\end{adjustbox}
\par\smallskip\parbox{0.95\linewidth}{\footnotesize \textit{Nota:} 13 economías, 2015-01 a 2026-07 (139 meses; 10,8 países por mes en promedio). Q1: Sharpe con la prueba de Ledoit y Wolf (2008), \textit{bootstrap} circular por bloques de 6 meses (4\,999 réplicas). Q3-H1: regresión cuantílica por momentos \citep{MachadoSantosSilva2019}, con efectos de país y de mes en la localización y de país en la escala; $p$ unilateral por \textit{bootstrap} de países (999). Q3-H2: regresión cuantílica recursiva con indicadoras de país, frente a un VaR normal con volatilidad exponencial ($\lambda=0{,}97$); prueba de Diebold y Mariano (1995) sobre la pérdida cuantílica, y de Christoffersen (1998). Desviaciones declaradas: el EWMA y el VaR histórico no se interrumpen en los meses sin dato; el VaR se evalúa donde existen los tres pronósticos; con umbrales expansivos la coincidencia ocurre en el 37\,\% de la muestra. Fuente: \texttt{p60\_senal\_dd.py} y \texttt{p62\_cola\_cambios.py}.}
\end{table}

\textbf{La fragilidad no ordena los retornos entre países} (Tabla~\ref{tab:quant-prereg}). Bruto, el Sharpe de la cartera es nulo y la pendiente de Fama--MacBeth no se distingue de cero; neto de costos es fuertemente negativo, porque una cartera mensual rota dos veces y media su tamaño. El resultado es el mismo sin Turquía, con costos de 25 puntos básicos, con pesos por rango, a tres meses y por subperíodos. La precedencia mensual es una regularidad dentro de cada país, no una señal de valor relativo entre países.

\textbf{La coincidencia tampoco ensancha la cola de los cambios del spread.} Cada riesgo por separado ensancha la distribución de los cambios, la fragilidad con evidencia firme. Su coincidencia no la ensancha más: el término conjunto es negativo, también a tres meses, sin 2020--2021, con los terciles de toda la muestra y sin efectos de mes. Es el mismo patrón subaditivo de la interacción en el nivel (Sección~\ref{sec:aditivos}). Un VaR que usa los dos riesgos no pronostica mejor que uno de volatilidad exponencial y no pasa la prueba de cobertura condicional. Según las consecuencias fijadas, la conclusión ``simultaneidad sin amplificación'' queda igual, ahora también para la cola de los cambios, y la precedencia mensual de la fragilidad no cambia de grado.

\section{Valor razonable, cobertura condicional y deuda en puntos básicos: pruebas preregistradas}
\label{sec:quant2-prereg}

Un octavo preregistro (\texttt{4\_Redaccion/preregistro\_quant2.md}), anterior a toda estimación, completa las preguntas de un portafolio de deuda soberana con las mismas bases mensuales, sin anticipación y desde 2015. Las tres hipótesis primarias son si los fundamentos de la tesis definen un valor razonable cuyo residuo pronostica el cambio del spread del mes siguiente (Q2), si una razón de cobertura frente al factor común de las demás economías mejora cuando depende de la fragilidad y de la deuda (Q4), y si la interacción entre la deuda y el riesgo de cola que predice el modelo de límite fiscal en puntos básicos (Sección~\ref{sec:modelo-resultados}) aparece con datos mensuales (Q5). Q5 usa un nowcast mensual del riesgo de cola (Q6), que debía validarse antes contra el $GaR$ oficial. Las ramas que exigían los CDS de los bancos o las fechas de las calificaciones soberanas se declararon no factibles con los datos del repositorio.

\begin{table}[htbp]
\centering
\small
\caption[Valor razonable, cobertura condicional y deuda en puntos básicos]{Valor razonable (Q2), cobertura condicional (Q4), nowcast del riesgo de cola (Q6) y deuda $\times$ riesgo de cola en puntos básicos (Q5): pruebas preregistradas}
\label{tab:quant2-prereg}
\begin{adjustbox}{max width=\linewidth}
\begin{tabular}{>{\raggedright\arraybackslash}p{4.6cm}>{\raggedright\arraybackslash}p{7.6cm}>{\raggedright\arraybackslash}p{2.6cm}}
\toprule
Prueba & Resultado & Holm \\
\midrule
Q2-H1: el residuo del valor razonable pronostica el cambio del spread a un mes & Clark--West $t=0{,}72$ ($p=0{,}24$); $R^2$ fuera de la muestra $+0{,}0006$; vida media del residuo 7,8 meses (criterio, menos de 6) & $0{,}47$: no pasa \\
\addlinespace
Q4-H1: cobertura condicional a la fragilidad y la deuda frente a un beta móvil de 36 meses & Diebold--Mariano $t=1{,}15$ ($p=0{,}13$); igual pérdida que un beta fijo; $c_{DD}=+0{,}025$ ($p=0{,}68$), $c_{deuda}=+0{,}16$ ($0{,}42$) & $0{,}37$: no pasa \\
\addlinespace
Q6: nowcast mensual del riesgo de cola (validación) & Correlación con el $GaR$ oficial $0{,}42$ (criterio, más de 0,9); error absoluto mediano 1,7 pp (criterio, menos de 0,5) & no valida \\
\addlinespace
Q5-H1: deuda $\times$ riesgo de cola en puntos básicos, mensual & $-1{,}1$ pb por d.e. de deuda (\textit{wild} unilateral $0{,}73$); positiva en 1 de 10 exclusiones; con el riesgo de cola oficial, $-0{,}9$ & $0{,}73$: no pasa (exploratoria) \\
\bottomrule
\end{tabular}
\end{adjustbox}
\par\smallskip\parbox{0.95\linewidth}{\footnotesize \textit{Nota:} 13 economías (10 con deuda trimestral en Q2 y Q5), 2015-01 a 2026-07. Q2: $\ln s$ sobre $\ln JLoss_{q-2}$, $D_{q-2}$ y la deuda de $q-4$ con efectos de país y de mes, reestimado cada mes; pronóstico $\hat\kappa\,\hat e$. Q4: factor, el cambio del logaritmo del promedio del spread de las otras doce economías; pérdida cuadrática del residuo cubierto. Q6: regresión cuantílica lineal ($\tau=0{,}05$) del crecimiento con el FCI y el VIX mensuales. Q5: spread en puntos básicos con efectos de país y de mes; \textit{wild cluster bootstrap} por país. Fuente: \texttt{p61}, \texttt{p63}, \texttt{p64} y \texttt{p65}.}
\end{table}

\textbf{Ninguna hipótesis pasa} (Tabla~\ref{tab:quant2-prereg}). El valor razonable no pronostica el cambio del spread del mes siguiente, y su residuo tarda casi ocho meses en cerrar la mitad de la brecha: la relación de largo plazo de la tesis describe niveles, pero no pronostica los cambios. La fragilidad y la deuda no cambian la exposición de cada país al factor común lo suficiente para mejorar la cobertura: un beta fijo por país cubre igual. El nowcast mensual del riesgo de cola no reproduce el $GaR$ oficial, que usa una distribución asimétrica ajustada y el crecimiento del trimestre, así que la prueba en puntos básicos queda como exploratoria; con el riesgo de cola oficial trimestral, la interacción en puntos básicos también es negativa. El modelo de límite fiscal sigue sin respaldo empírico en puntos básicos.
